\originalchapter{8.05 2011期末：氢、合成自旋与压缩态}
原件为Fall 2011 Final Exam, 2011年12月22日9:00--12:00, 七题100分, 闭卷禁计算器; Fall 2013 OCW归档. 以下为\textbf{AI独立解答}, 保留原卷参数歧义并说明其影响.
\originalsection{第1题：十个判断（20分）}
\begin{exercise}
判断: (1)Hermitian算符反对易子Hermitian; (2)各时刻Schrödinger $H$对易则$H_H=H_S$; (3)$a_0\alpha$为经典电子半径; (4)$a_0\alpha^2$为电子Compton波长; (5)因子分解$H^{(1)},H^{(2)}$至少一个必有正规零能态; (6)两个独立角动量之差也满足角动量代数; (7)$\boldsymbol J_1\cdot\boldsymbol J_2$与$J_{1x}+J_{2x}$对易; (8)任意自旋$j$表示各$J_i$迹0; (9)三维球对称振子第一激发组是$\ell=1$多重态; (10)第二激发组全是$\ell=2$多重态.
\end{exercise}
\begin{solution}
依次\textbf{T,T,F,F,F,F,T,T,T,F}. (1)$(AB+BA)^\dagger=AB+BA$. (2)$U$与当时$H$对易, 共轭不变. (3)$a_0\alpha$是约化Compton长度, (4)$a_0\alpha^2$才是经典半径, 非约化Compton另有$2\pi$. (5)$A\psi=0$的形式解可能不可归一, 例如全线常数超势$W=1$给指数向一侧发散, 两Hamiltonian都没有正规零能态. (6)差的自对易右侧给$J_1+J_2$, 不是差. (7)它是共同旋转标量. (8)$J_z$谱$m=-j,\ldots,j$相加0, 其他方向等谱. (9)三维第一组维数3, 三个单量子态是向量表示. (10)第二组六维分成$\ell=2$五维加$\ell=0$一维.
\end{solution}
\originalsection{第2题：Heisenberg升降算符（10分）}
\begin{exercise}
求谐振子的$a_H(t),a_H^\dagger(t),N_H(t)$.
\end{exercise}
\begin{solution}
$[H,a]=-\hbar\omega a$, 因而$\dot a_H=-\ii\omega a_H$, 解$a_H=\e^{-\ii\omega t}a$, $a_H^\dagger=\e^{\ii\omega t}a^\dagger$. 两相位抵消,$N_H=N$. 能量占有数不变, 相位旋转使位置动量均值和某些宽度仍可随时间变化.
\end{solution}
\originalsection{第3题：倒数均值不等式（15分）}
\begin{exercise}
(a)从$\langle u^{-1}(u-\langle u\rangle)^2\rangle\geq0$推出$\langle1/u\rangle$与$1/\langle u\rangle$关系,$u$谱非负. (b)以氢基态$r$验证, 可用维里定理求$\langle1/r\rangle$.
\end{exercise}
\begin{solution}
需$u^{-1}$在所取态有意义, $\langle u\rangle>0$, 有关矩有限. 展开给$\langle u\rangle^2\langle1/u\rangle-\langle u\rangle\geq0$, 即$\langle1/u\rangle\geq1/\langle u\rangle$. 零谱成分可使逆不存在或均值无穷, 不能直接除0. 氢基态径向概率$4r^2\e^{-2r/a_0}/a_0^3$给$\langle r\rangle=3a_0/2$. 维里$2T=-V$与$E=-e^2/(2a_0)$给$\langle V\rangle=2E=-e^2/a_0$, 所以$\langle1/r\rangle=1/a_0>2/(3a_0)$. 位置有分散, 故严格不等号.
\end{solution}
\originalsection{第4题：氢Hamiltonian的正性分解（15分）}
\begin{exercise}
将$H=\boldsymbol p^2/(2m)-e^2/r$写为$\sum_i(p_i+\ii\beta x_i/r)(p_i-\ii\beta x_i/r)/(2m)+\gamma$,$\beta,\gamma$实. (a)求常数. (b)证所有态$\langle H\rangle\geq\gamma$. (c)求取等号的径向波函数.
\end{exercise}
\begin{solution}
令$n_i=x_i/r$, $\sum n_i^2=1$, $\sum\partial_i n_i=2/r$（在原点外）. 展开乘积得到$p^2+\beta^2-2\beta\hbar/r$, 因而$\beta=me^2/\hbar$, $\gamma=-me^4/(2\hbar^2)$. 每项是$Q_i^\dagger Q_i$, $Q_i=p_i-\ii\beta n_i$, 所以均值为$\sum\norm{Q_i\psi}^2/(2m)+\gamma\geq\gamma$, 正规Coulomb域允许这一二次型关系. 等号要求全部$Q_i\psi=0$, 对径向$\psi$即$\hbar\psi'(r)+\beta\psi=0$, 得$\psi=(\pi a_0^3)^{-1/2}\e^{-r/a_0}$. 这既给真实基态又证明能量下界恰可取到, 不是只猜出指数试探态.
\end{solution}
\originalsection{第5题：三个自旋1的反铁磁圆环（15分）}
\begin{exercise}
三可区分自旋1有$H=\Delta(\boldsymbol S_1\cdot\boldsymbol S_2+\boldsymbol S_2\cdot\boldsymbol S_3+\boldsymbol S_3\cdot\boldsymbol S_1)/\hbar^2$,$\Delta>0$. (a)维数并改写$H$. (b)求各能量简并度. (c)在$\ket{m_1,m_2,m_3}$基中求归一基态, 从$m_1+m_2+m_3=0$的七候选态开始.
\end{exercise}
\begin{solution}
维数$3^3=27$, 总$\boldsymbol S=\sum\boldsymbol S_i$给$H=\Delta[S^2-6\hbar^2]/(2\hbar^2)$. 合成$1\otimes1\otimes1=0\oplus3(1)\oplus2(2)\oplus3$, 这里数字前系数为副本数, 最后3是$j=3$不是三个副本. 能量和总简并为
\begin{center}\begin{tabular}{cccc}\toprule $j$&副本数&$E/\Delta$&简并度\\\midrule0&1&$-3$&1\\1&3&$-2$&9\\2&2&0&10\\3&1&3&7\\\bottomrule\end{tabular}\end{center}
总和27. 唯一$j=0$态需$S_z\psi=S_+\psi=0$. 令000的系数为$g$, 六排列$+0-,+-0,0+-,0-+,-+0,-0+$的系数依次$a,b,c,d,e,f$. 除去共同升阶因子$\sqrt2\hbar$, $S_+\psi=0$在$++-,+-+,-++$上的约束为
\[
a+c=0,\qquad b+d=0,\qquad e+f=0.
\]
在$+00,0+0,00+$上的约束为$a+b+g=0,c+e+g=0,d+f+g=0$. 后三式相加并用前三式得$3g=0$; 再逐项解得$(a,b,c,d,e,f)=a(1,-1,-1,1,1,-1)$. 归一结果为
\[
\ket{\psi_g}=\frac1{\sqrt6}\bigl(\ket{+,0,-}+\ket{0,-,+}+\ket{-,+,0}
-\ket{+,-,0}-\ket{-,0,+}-\ket{0,+,-}\bigr).
\]
可逐项作用$S_+$, 每个目标排列出现一正一负, 完全抵消. 唯一副本数保证除整体相位无其他基态.
\end{solution}
\originalsection{第6题：因子分解的阶梯结构（10分）}
\begin{exercise}
假设$H_\ell=H_\ell^{(1)}+\alpha_\ell$, $H_\ell^{(1)}=A_\ell^\dagger A_\ell$, $H_\ell^{(2)}=A_\ell A_\ell^\dagger=H_{\ell+1}^{(1)}+\beta_\ell$. 每个$H_\ell^{(1)}$有已知正规零态$\phi_{0,\ell}^{(1)}$. (a)写$H_\ell$前三态及能量. (b)推广任意$n$.
\end{exercise}
\begin{solution}
恒等$H_\ell^{(1)}A_\ell^\dagger=A_\ell^\dagger H_\ell^{(2)}$把伙伴能量搬回. 因此$\phi_{0,\ell}\propto\phi_{0,\ell}^{(1)}$, $E_{0,\ell}=\alpha_\ell$; $\phi_{1,\ell}\propto A_\ell^\dagger\phi_{0,\ell+1}^{(1)}$, $E_{1,\ell}=\alpha_\ell+\beta_\ell$; $\phi_{2,\ell}\propto A_\ell^\dagger A_{\ell+1}^\dagger\phi_{0,\ell+2}^{(1)}$, $E_{2,\ell}=\alpha_\ell+\beta_\ell+\beta_{\ell+1}$. 一般
\[
\phi_{n,\ell}\propto A_\ell^\dagger A_{\ell+1}^\dagger\cdots A_{\ell+n-1}^\dagger\phi_{0,\ell+n}^{(1)},\quad
E_{n,\ell}=\alpha_\ell+\sum_{q=0}^{n-1}\beta_{\ell+q}.
\]
归一因子由每步范数平方等于所提升的伙伴能量求得. 构造需输出非零、正规且满足同一边界域; 题给链结构可产生这些态, 若要声称某任意超势的\textbf{全部}谱还须排除链外额外态.
\end{solution}
\originalsection{第7题：压缩真空的时间宽度（15分）}
\begin{exercise}
题面$S(\gamma)=\exp[-\gamma(a^{\dagger2}-a^2)/2]$, $\ket{0_\gamma}=S(\gamma)\ket0$. (a)求$S^\dagger xS,S^\dagger pS$. (b)求演化后$(\Delta x(t))^2=\hbar G/(2m\omega)$. 原题此处印$\gamma\in\mathbb C$, 但同卷公式表第5页限定$\gamma\in\mathbb R$.
\end{exercise}
\begin{solution}
必须先处理条件: 题给生成元$a^{\dagger2}-a^2$为反Hermitian, 只有\textbf{实}$\gamma$才使$S$幺正, 保证题中“压缩真空”的单位归一. 按公式表一致的实参数解, 用$S^\dagger aS=a\cosh\gamma-a^\dagger\sinh\gamma$, 得$S^\dagger xS=\e^{-\gamma}x$, $S^\dagger pS=\e^\gamma p$. Heisenberg$x_H=x\cos\omega t+p\sin\omega t/(m\omega)$, 真空交叉对称矩0, 所以
\[
G(\gamma,\omega t)=\e^{-2\gamma}\cos^2\omega t+\e^{2\gamma}\sin^2\omega t.
\]
低噪声方向随振子旋转, 能量均值仍不变.

若坚持题面任意复$\gamma$, $S^\dagger S\neq I$, 前述共轭缩放公式和归一化不能成立, 有时作用真空甚至不正规. 可在$|w|<1$, $w=-\tanh\gamma$时先将形式态正规化, 其位置宽度为$G=[1+|w|^2+2\operatorname{Re}(w\e^{-2\ii\omega t})]/(1-|w|^2)$. 这不是用原题非幺正$S^\dagger xS$直接得出的单位态公式. 通常复压缩参数应改生成元为$(\zeta^*a^2-\zeta a^{\dagger2})/2$; 本册把它列为条件修正说明, 不伪称原卷已印此式.
\end{solution}
